% Part: sets-functions-relations
% Chapter: relations
% Section: equivalence-relations

\documentclass[../../../include/open-logic-section]{subfiles}

\begin{document}

\olfileid{sfr}{rel}{eqv}

\olsection{સામ્ય સંબંધો}

ગણ પરનો તાદાત્મ્ય સંબંધ સ્વવાચક, સંમિત અને પરંપરિત છે.
આ ત્રણેય ગુણધર્મો ધરાવતા સંબંધો $R$ ઘણા સામાન્ય છે.

\begin{defn}[સામ્ય સંબંધ]
સ્વવાચક, સંમિત અને પરંપરિત એવા સંબંધ $R \subseteq A^2$ને
\emph{સામ્ય સંબંધ} કહે છે. $A$ના !!^{element}s $x$ અને $y$
એ \emph{$R$-સામ્ય} ધરાવે છે એમ કહેવાય જો $Rxy$.
\end{defn}

સામ્ય સંબંધોથી \emph{સામ્ય વર્ગ}નો ખ્યાલ મળે છે.
સામ્ય સંબંધ પ્રદેશને જુદા જુદા વર્ગોમાં “વહેંચે છે”.
દરેક વર્ગની બધી વસ્તુઓ એકબીજા સાથે સંબંધ ધરાવે છે;
અને જુદા વર્ગોની કોઈ વસ્તુઓ એકબીજા સાથે સંબંધ ધરાવતી નથી.
ક્યારેક આ વર્ગોની જ \emph{સીધી} વાત કરવી ઉપયોગી બને છે.
તે માટે આપણે એક વ્યાખ્યા રજૂ કરીએ:

\begin{defn}\ollabel{def:equivalenceclass}
ધારો કે $R \subseteq A^2$ સામ્ય સંબંધ છે. દરેક $x \in A$
માટે, $A$માં $x$નો \emph{સામ્ય વર્ગ} એ ગણ
$\equivrep{x}{R} = \Setabs{y \in A}{Rxy}$ છે.
$R$ અનુસાર $A$નો \emph{ભાગફળ ગણ} એ
$\equivclass{A}{R} = \Setabs{\equivrep{x}{R}}{x \in A}$ છે,
એટલે કે આ સામ્ય વર્ગોનો ગણ.
\end{defn}

સામ્ય વર્ગો ખરેખર $A$નું વર્ગ વિભાજન કરે છે તે સાબિત કરીને,
હવે પછીનું પરિણામ સામ્ય વર્ગની વ્યાખ્યાને સમર્થન આપે છે:

\begin{prop}
જો $R \subseteq A^2$ સામ્ય સંબંધ હોય, તો $Rxy$ હોય
તો અને તો જ $\equivrep{x}{R} = \equivrep{y}{R}$.
\end{prop}

\begin{proof}
ડાબેથી જમણી દિશા માટે, ધારો કે $Rxy$, અને
$z \in \equivrep{x}{R}$ લો. તો વ્યાખ્યા પ્રમાણે $Rxz$.
$R$ સામ્ય સંબંધ હોવાથી, $Ryz$. (વિગતે કહીએ તો:
$Rxy$ અને $R$ સંમિત હોવાથી $Ryx$ મળે છે, અને
$Rxz$ અને $R$ પરંપરિત હોવાથી $Ryz$ મળે છે.)
તેથી $z \in \equivrep{y}{R}$. આથી સામાન્ય રીતે
$\equivrep{x}{R} \subseteq \equivrep{y}{R}$.
પણ બરાબર એ જ રીતે,
$\equivrep{y}{R} \subseteq \equivrep{x}{R}$.
તેથી ઘટકો દ્વારા નિર્ધારિત સમાનતાના સિદ્ધાંતથી
$\equivrep{x}{R} = \equivrep{y}{R}$.

જમણેથી ડાબી દિશા માટે, ધારો કે
$\equivrep{x}{R} = \equivrep{y}{R}$.
$R$ સ્વવાચક હોવાથી $Ryy$, તેથી $y \in \equivrep{y}{R}$.
આમ $\equivrep{x}{R} = \equivrep{y}{R}$ એવી ધારણાથી
$y \in \equivrep{x}{R}$ પણ મળે છે. તેથી $Rxy$.
\end{proof}

\begin{ex}
મોડ્યુલો અંકગણિતમાંથી સામ્ય સંબંધોનો સુંદર દાખલો મળે છે.
કોઈ પણ $a$, $b$ અને $n \in \PosInt$ માટે,
$a \equiv_n b$ એમ કહીએ જો અને ફક્ત જો $a$ને $n$ વડે
ભાગતાં જે શેષ મળે તે જ શેષ $b$ને $n$ વડે ભાગતાં મળે.
(થોડું વધુ સાંકેતિક રીતે: $a \equiv_n b$ તો અને તો જ,
કોઈક $k \in \Int$ માટે, $a - b = kn$.)
હવે, કોઈ પણ $n$ માટે, $\equiv_n$ સામ્ય સંબંધ છે.
અને $\equiv_n$થી બરાબર $n$ જુદા સામ્ય વર્ગો મળે છે;
એટલે કે, $\equivclass{\Nat}{\equiv_n}$માં $n$ !!{element}s છે.
આ વર્ગો છે: $n$ વડે નિઃશેષ ભાગી શકાય તેવી સંખ્યાઓનો ગણ,
એટલે કે $\equivrep{0}{\equiv_n}$; $n$ વડે ભાગતાં
શેષ $1$ મળે તેવી સંખ્યાઓનો ગણ, એટલે કે
$\equivrep{1}{\equiv_n}$; \ldots; અને $n$ વડે ભાગતાં
શેષ $n-1$ મળે તેવી સંખ્યાઓનો ગણ, એટલે કે
$\equivrep{n-1}{\equiv_n}$.
\end{ex}

\begin{prob}
દરેક $n \in \PosInt$ માટે $\equiv_n$ સામ્ય સંબંધ છે અને
$\equivclass{\Nat}{\equiv_n}$માં બરાબર $n$ સભ્યો છે તે દર્શાવો.
\end{prob}

\end{document}
